\section{Tick size}

\begin{qfdefinition}[Tick size: the price grid, tick regimes and relative tick, when the tick binds, how it prices queue jumping, and price clustering]\label{def:tick-size}
\qffield{Statement} The tick size defines the price grid on which a limit order must lie: every valid price is a multiple of it, so it is also the smallest difference between two limit prices \citep[chap.~1, sec.~1.3.1, p.~74]{alpha-laruelle}. It may apply to every stock whatever its price and liquidity (as in the United States for stocks priced above \$1) or depend on price and liquidity bands, a \emph{tick-size regime} (p.~74). The \emph{relative tick} $\vartheta/m(t)$, tick over mid-price, compares instruments: Priceline at about \$1{,}000 in 2014 had a relative tick near $10^{-5}$ with the same one-cent tick as a \$1 stock \citep[chap.~3, sec.~3.1.5, pp.~48--49]{catalogue-0528}. A \emph{large-tick} instrument has a spread close to one tick most of the time, a \emph{small-tick} one a spread of many ticks (chap.~4, p.~58).

\qffield{Conditions} Rules and figures as of 2002--2015; U.S.\ and EU tick regimes have changed since and are not checked here. Being the smallest price increment, the tick is a floor on the spread, so a smaller tick lowers spreads only where that floor binds \citep[chap.~1, sec.~1.3.2, p.~77]{alpha-laruelle}. Two simple measures of how much it binds (p.~77): the \emph{binding probability}, the share of trades, value traded or quote duration with a one-tick spread, and the \emph{spread leeway}, the spread in ticks minus one. For S\&P~500 stocks in July 2012 most stocks below \$20 had a one-tick spread most of the time, and stocks above \$150 showed no constraint (Figures~1.19--1.20, pp.~78--79).

\qffield{Practice} The tick sets the price of priority: where it is very small, a trader can step ahead of resting orders by one tick for little cost, which leaves little volume at any level and makes the best quotes change very often \citep[chap.~3, sec.~3.1.5, p.~49]{catalogue-0528}; Lehalle and Laruelle add that queue jumping with a small tick ``makes time priority nearly meaningless'' and quotes more fleeting, while with a large tick and a one-tick spread an impatient trader must cross and a liquidity provider must manage its place in the queue \citep[chap.~1, pp.~88--89]{alpha-laruelle}.

Harris reports that prices cluster on round numbers (whole numbers more than halves, halves more than odd quarters; in decimal markets multiples of 1.00, 0.50, 0.25, 0.20, 0.10 and 0.05), most when the minimum increment is a small fraction of price, the market volatile and the instrument thinly traded, and that traders place orders just above or below round numbers to step ahead of others \citep[chap.~5, box ``Price Clustering'', p.~91]{alpha-harris}.

\emph{Ticks, precedence and competition between venues}. Time precedence ``is meaningful only when the minimum price increment is not trivially small'': too small a tick weakens precedence and price competition, too large a tick makes improving the price expensive \citep[chap.~6, sec.~6.1.1.2, p.~114]{alpha-harris}. Large traders who expose orders prefer strict time precedence with an economically significant tick; small market-order traders prefer small ticks, since the spread cannot be smaller than the tick \citep[chap.~26, sec.~26.3.1, p.~531]{alpha-harris}. A market cannot keep a meaningful tick when competitors trading the same instrument offer a smaller one: small markets with few limit orders cut ticks to attract quote matchers, and larger ones must follow; Harris's sidebar \emph{The Race to the Bottom} reads the U.S.\ tick reductions of the 1990s, which ``culminated in the full decimalization of the markets in 2001'', this way (sec.~26.5.2, pp.~537--538).

In Europe, Lehalle and Laruelle date a ``tick war'' to June 2009, when BATS and Turquoise cut tick sizes, ``leading to a tacit agreement between exchanges under the FESE'' (Table~1.1, p.~15); MiFID~II then required regulated markets to adopt tick-size regimes calibrated to liquidity and average spread (Art.~49, \hyperref[rule:eu-trading-venue-transparency]{EU trading venues and transparency}). The current U.S.\ increment rule and the stayed half-penny tick are in \hyperref[rule:us-equity-market-structure]{U.S.\ equity market structure rules}. The ``who gains from tick changes'' list and its tick-policy principle (vary ticks with price and liquidity; revise them on spreads, depth, order-to-trade ratios and standard-size costs) are a synthesis of these arguments, not a rule.

\qffield{Pitfalls} Explaining price clustering by communication, search and queueing convenience requires evidence; Harris states when clustering is strongest, not why, and those reasons are not adopted. Their price- and liquidity-banded regimes match Lehalle and Laruelle's price ranges and quotation groups, which ``theoretically'' group stocks by liquidity (p.~74). Their claim that once the tick stops binding ``further reductions mostly change queue behaviour'' goes beyond the passages: Lehalle and Laruelle report spreads that did not fall, and sometimes rose, for illiquid high-priced stocks after decimalisation (\hyperref[ev:tick-size-market-quality]{tick size and market quality}).

\qffield{See also} \hyperref[mod:harris-discretised-spread]{Harris's discretised spread}; \hyperref[ev:tick-size-market-quality]{tick size and market quality}; \hyperref[res:spread-determinants]{spread determinants} (large-tick bounce and depth); \hyperref[def:queue-imbalance]{queue imbalance} (large-tick races); \hyperref[mech:order-lifecycle]{order lifecycle} (priority rules).

\qffield{Further reading} The pathway's readings on the U.S.\ tick-size pilot and the 2024 tick and access-fee rule (Barardehi et al.; the SEC release) are not read here, so no current U.S.\ tick schedule is stated.

\end{qfdefinition}

\begin{qfevidence}[Tick size and market quality: quoted depth, order and trade sizes, order-to-trade ratios, cumulative depth and the evidence from tick reductions]\label{ev:tick-size-market-quality}
\qffield{Statement} On FTSE~100 stocks in April 2012, grouped by deciles of the daily spread in ticks, Lehalle and Laruelle find far less volume at the best limits when the tick is small relative to the spread, a nearly threefold fall in the size of orders at the first limit from the first to the last decile, smaller trades, and a 50 percent rise in the number of book updates per trade \citep[chap.~1, sec.~1.3.2, Figures~1.21--1.24, pp.~80--84]{alpha-laruelle}. They read the last as a fall in the value of a quote, smaller and shorter-lived (p.~84).

From the decimalisation studies they cite: spreads fell for active low-priced stocks whose spread had been one tick, with ``at best no effect'' and some increases for illiquid high-priced stocks, and a tick below which further reductions widen spreads (Ahn et al.\ found the direction of Harris's prediction right but its magnitude overestimated) (pp.~79--80). Harris's prediction was a cross-sectional model projection, not an observed change: for a \$1/16 tick he projected, for stocks priced under \$10, spreads 38 percent narrower, quotation sizes 16 percent smaller and daily volume 34 percent higher \citep[p.~149]{harris-1994}.

When the NYSE moved from \$1/8 to \$1/16 on 24 June 1997, Goldstein and Kavajecz found, for 100 common stocks from the TORQ sample (plus 10 funds) over one four-week period before and two after, quoted spreads \$0.03 (14.3 percent) narrower and quoted depth 48 percent lower, while the spread of the limit-order book widened by \$0.03 (9.1 percent), its depth at the best prices fell 48 percent and cumulative depth fell throughout the book; small orders gained, large orders lost, and the gains concentrated in frequently traded stocks \citep[abstract and results, PDF pp.~2, 5 and~9--15]{goldstein-kavajecz-1998}.

\qffield{Conditions} FTSE~100 data of April 2012; decimalisation studies of 1992--2005. Spreads alone do not measure quality: a smaller quoted spread may be available for less volume (p.~81). \emph{Cumulative depth} rebuilt from order-level data measures the cost of a marketable order of any size and includes hidden liquidity, but conditions on how liquidity takers changed their order sizes; Bacidore et al.\ found NYSE marketable orders cheaper at every size after decimalisation and depth thinner beyond 15 cents from the best quotes, and Goldstein and Kavajecz, with costs in basis points, found the public book less deep and large trades costlier (pp.~85--87). In a fragmented market a single book is not the market, and mirrored orders make aggregation across venues unreliable (pp.~75--76, 85).

These are cross-sectional associations (FTSE~100) and before--after comparisons around single events. Goldstein and Kavajecz's sample is the surviving TORQ stocks, slightly overweight in large stocks (n.~7, PDF p.~9), and the read version is their September 1998 working paper, not the 2000 journal article. Harris's figures are projections from his 1994 model, whose second sample (second quarter of 1989, tickers EN to MOT, about a quarter of NYSE and AMEX firms) had 48 percent of quotation spreads at \$1/8 on average and a third even for stocks priced over \$40 (p.~155); his projection tables are in the part of the scan that does not render (PDF pp.~27--30). The other studies are reported by Lehalle and Laruelle and were not read.

\qffield{Practice} The proxies for tick-size effects, binding probability, spread leeway, depth near the mid-price and message intensity (order-to-trade ratio, cancellations and modifications), are the measures Lehalle and Laruelle use (pp.~77 and~84); see \hyperref[def:tick-size]{tick size}. More messages per trade need not mean more liquidity.

\qffield{Pitfalls} The caveat that ``a consolidated tape of trades cannot fully replace quote data if only the best limit is visible'' was not found in these passages and is not adopted; Lehalle and Laruelle's point is the falling value of each quote. The ``cumulative depth'' mechanism omitted the conditioning on takers' behaviour and the fragmentation problem.

\qffield{See also} \hyperref[def:tick-size]{tick size}; \hyperref[mod:harris-discretised-spread]{Harris's discretised spread}; \hyperref[def:spread-measures]{spread measures}; \hyperref[def:best-execution]{best execution}; \hyperref[def:smart-order-routing]{smart order routing} (fragmented books).

\qffield{Further reading} \citep{goldstein-kavajecz-1998} (working paper; the 2000 journal version not read) and \citep{harris-1994} (partial scan). Bacidore et al.\ (2003), Jones and Lipson (2001), Ahn et al.\ (1996), Harris (1996) and Aitken and Comerton-Forde (2005), cited by Lehalle and Laruelle, are not read.

\end{qfevidence}

\begin{qfmodel}[Harris's discretised spread: a latent continuous spread rounded to the tick, its fit by maximum likelihood, and the counterfactual spread under a smaller tick]\label{mod:harris-discretised-spread}
\qffield{Statement} Introduced by Harris (1994) to forecast the effect of decimalisation, as presented by Lehalle and Laruelle \citep[app.~A, sec.~A.6, pp.~266--268]{alpha-laruelle}: for stock $j$, trade $t$ on day $d$, the observed spread in ticks is $\Psi_{j,t}=(\mathrm{Ask}_{j,t}-\mathrm{Bid}_{j,t})/\mathrm{tick}_{j,d}$, using only days with a single tick, and $\Psi_{j,t}=\Phi(\tilde\Psi_{j,t})$ for a latent continuous spread $\tilde\Psi_{j,t}$, where $\Phi$ gives the nearest integer for inputs above 1.5 and 1 otherwise. The latent spread has density $\tilde f_{s,\theta}$ with mean
\[ E[\tilde\Psi_{j,t}\mid T_j,\sigma_{j,d}]=g(T_j,\sigma_{j,d})\,\frac{P_{j,d}}{\mathrm{tick}_{j,d}},\qquad g=\alpha\,T_j^{\beta}\sigma_{j,d}^{\gamma}, \]
with $T_j$ mean turnover, $\sigma_{j,d}$ volatility and $P_{j,d}$ the mean traded price; $g$ is the latent relative spread and $P/\mathrm{tick}$ converts it to ticks (pp.~268--269). The discrete probabilities are $f_{s,\theta}(1)=\int_0^{1.5}\tilde f_{s,\theta}$ and $f_{s,\theta}(i)=\int_{i-0.5}^{i+0.5}\tilde f_{s,\theta}$ for $i\ge2$, and $(\alpha,\beta,\gamma,\theta)$ maximise $\sum_{j,d,k}N(j,d,k)\ln f(k)$ over the counts $N(j,d,k)$ of $k$-tick spreads (pp.~269--270).

\qffield{Notation} Counterfactual: for a new tick of old tick$/\delta$, the latent spread in \emph{new} ticks is $\delta\tilde\Psi$, rounded by the same rule, so $f_{s,\theta,\delta}(1)=\int_0^{1.5/\delta}\tilde f_{s,\theta}$ and $f_{s,\theta,\delta}(k)=\int_{(k-0.5)/\delta}^{(k+0.5)/\delta}\tilde f_{s,\theta}$ for $k\ge2$, and the expected spread in \emph{old} ticks is
\[ E(\Psi_{j,d}\mid T_j,\sigma_{j,d},\delta)=\sum_k\frac k\delta\,f_{\hat\alpha T_j^{\hat\beta}\sigma_{j,d}^{\hat\gamma}P_{j,d}/\mathrm{tick}_{j,d},\hat\theta,\delta}(k) \]
(pp.~272--273, checked by page image): $k$ counts new ticks and $k/\delta$ converts them to old ones. An observed spread of $k$ old ticks is $k\delta$ new ticks.

\qffield{Conditions} The latent distribution and its dependence on turnover and volatility are held fixed when the tick changes; that is the model's counterfactual assumption, which ignores changes in order exposure, depth and participants (\hyperref[ev:tick-size-market-quality]{tick size and market quality}). Lehalle and Laruelle fit a gamma latent distribution to the 60 most liquid FTSE stocks in April 2012, with standardised $\beta'=-0.46$ and $\gamma'=2.62$ (more liquidity, lower spread; more volatility, higher), and a pseudo-$R^2$ of 0.8 between observed and expected relative spreads (pp.~270--271).

\qffield{Practice} The expected decrease is largest for stocks whose spread is near one tick (Figure~A.9, p.~272). Checked (\nolinkurl{order-flow-checks.py}), with a gamma latent spread of shape 4 and $\delta=5$: a latent mean of 0.6 old ticks gives an expected observed spread of 1.010 old ticks now and 0.601 after the change ($-41$ percent); a mean of 1.2 gives 1.302 and 1.200 ($-8$ percent); means of 3 and 8 change by less than 0.01 old ticks.

\qffield{Pitfalls} The counterfactual said ``an observed $k$-tick spread becomes roughly $k/\delta$ old-tick units before re-rounding''. That confuses the units: an observed $k$-tick spread is $k\delta$ new ticks, and $k/\delta$ old ticks is the value of $k$ \emph{new} ticks; nor does the model rescale observed spreads, it re-rounds the latent spread (for a latent mean of 8 old ticks the expected change is nil, not $8/5-8$). Lehalle and Laruelle's caption to Figure~A.10 gives the change ``in old ticks unit'' while the text calls it a relative decrease (p.~273); the figure's axis is a difference in old ticks. Their Figure~A.8 fit is in-sample.

\qffield{See also} \hyperref[def:tick-size]{tick size}; \hyperref[ev:tick-size-market-quality]{tick size and market quality}; \hyperref[res:spread-determinants]{spread determinants}; \hyperref[meth:spread-volatility-regression]{spread--volatility regression}.

\qffield{Further reading} Harris (1994), cited by Lehalle and Laruelle (p.~74), is not read. Dayri and Rosenbaum's paper on large-tick assets (arXiv:1207.6325v2) provides futures evidence (\hyperref[res:spread-determinants]{spread determinants}); its uncertainty-zone model is not summarised here.

\end{qfmodel}
